\newcommand{\StudyBook}{CS5222-Final}
\newcommand{\StudyRoot}{../../..}
\documentclass[11pt,lang=cn,chinesefont=nofont,mode=simple,math=newtx,color=blue]{elegantbook}
\usepackage{longtable,array,fvextra,xurl,adjustbox,needspace}
\usepackage[normalem]{ulem}
\providecommand{\lt}{<}
\providecommand{\gt}{>}
\tracinglostchars=3
\setCJKmainfont{FandolSong-Regular.otf}[BoldFont=FandolSong-Bold.otf,ItalicFont=FandolKai-Regular.otf]
\setCJKsansfont{FandolHei-Regular.otf}[BoldFont=FandolHei-Bold.otf]
\setCJKmonofont{FandolFang-Regular.otf}
\setmonofont{DejaVuSansMono.ttf}[BoldFont=DejaVuSansMono-Bold.ttf,ItalicFont=DejaVuSansMono-Oblique.ttf,Scale=.84]
\geometry{a4paper,left=22mm,right=18mm,top=20mm,bottom=21mm,headheight=15pt,headsep=7mm}
\definecolor{structurecolor}{HTML}{173F57}
\colorlet{main}{structurecolor}
\colorlet{second}{structurecolor}
\colorlet{third}{structurecolor}
\hypersetup{unicode=true,colorlinks=true,linkcolor=structurecolor,urlcolor=structurecolor,filecolor=structurecolor,citecolor=structurecolor,pdfauthor={CityU StudyNotes}}
\makeatletter
\renewcommand{\tableofcontents}{%
  \hypersetup{linktoc=all,linkcolor=black}%
  \chapter*{\contentsname}%
  \StudyTarget{\StudyBook__contents}\StudyTarget{\StudyBook___1}%
  \@starttoc{toc}%
  \hypersetup{linkcolor=structurecolor}\clearpage}
\makeatother
\setstretch{1.25}
\setlength{\parskip}{.35\baselineskip}
\setlength{\parindent}{0pt}
\setlength{\emergencystretch}{2em}
\setcounter{tocdepth}{1}
\setcounter{secnumdepth}{-1}
\setlist{nosep,leftmargin=1.8em,topsep=.35\baselineskip}
\titleformat{\chapter}[display]{\Huge\sffamily\bfseries\color{structurecolor}}{}{0pt}{}[\vspace{.4ex}\titlerule]
\titlespacing*{\chapter}{0pt}{0pt}{22pt}
\titleformat{\section}{\Large\sffamily\bfseries\color{structurecolor}}{}{0pt}{}
\titleformat{\subsection}{\large\sffamily\bfseries}{}{0pt}{}
\titleformat{\subsubsection}{\normalsize\sffamily\bfseries}{}{0pt}{}
\titlespacing*{\section}{0pt}{1.6\baselineskip}{.5\baselineskip}
\titlespacing*{\subsection}{0pt}{1.1\baselineskip}{.35\baselineskip}
\providecommand{\StudyBook}{Sample}
\providecommand{\StudyRoot}{../..}
\providecommand{\StudySourceRef}{main}
\providecommand{\StudyDate}{2026-10-16}
\newcommand{\StudyRunning}{\StudyBook}
\fancyhf{}
\fancyhead[L]{\small\sffamily\color{gray}\StudyRunning}
\fancyhead[R]{\small\sffamily\color{gray}CityU StudyNotes}
\fancyfoot[R]{\small\thepage}
\renewcommand{\headrulewidth}{.3pt}
\pagestyle{fancy}
\fancypagestyle{plain}{\fancyhf{}\fancyfoot[R]{\small\thepage}\renewcommand{\headrulewidth}{0pt}}
\newwrite\StudyMap
\AtBeginDocument{\immediate\openout\StudyMap=\jobname.pages}
\makeatletter
\newcommand{\StudyTarget}[1]{%
  \hypertarget{#1}{}\label{#1}%
  \protected@write\StudyMap{}{#1|\thepage}%
}
\makeatother
\newcommand{\StudyChapter}[2]{%
  \clearpage\renewcommand{\StudyRunning}{\ifcsname studytitle@#1\endcsname\csname studytitle@#1\endcsname\else#1\fi}%
  \chapter*{#2}\phantomsection\addcontentsline{toc}{chapter}{#2}%
  \StudyTarget{#1__root}%
}
\newcommand{\StudyHeading}[2]{%
  \ifcase#1\relax\or
    \section*{#2}\phantomsection\addcontentsline{toc}{section}{#2}%
  \or\subsection*{#2}\or\subsubsection*{#2}\else\paragraph*{#2}\fi
}
\newcommand{\StudyImage}[1]{%
  \par\medskip\begingroup
  \edef\StudyResolvedImage{\ifcsname studyimage@#1\endcsname\csname studyimage@#1\endcsname\else\StudyRoot/#1\fi}%
  \noindent\makebox[\linewidth][c]{%
  \adjustbox{max width=.96\linewidth,max height=.69\textheight}{%
  \includegraphics{\StudyResolvedImage}}}\endgroup\par\medskip
}
\newcommand{\StudyCellImage}[1]{%
  \begingroup
  \edef\StudyResolvedImage{\ifcsname studyimage@#1\endcsname\csname studyimage@#1\endcsname\else\StudyRoot/#1\fi}%
  \adjustbox{max width=\linewidth,max height=.28\textheight}{\includegraphics{\StudyResolvedImage}}\endgroup}
\DefineVerbatimEnvironment{StudyCode}{Verbatim}{breaklines=true,breakanywhere=true,fontsize=\small,tabsize=4,framesep=3mm,frame=leftline,rulecolor=\color{gray!55}}
\newcommand{\StudyInline}[1]{\allowbreak{}\texttt{#1}\allowbreak{}}
\newenvironment{StudySolution}[1]{%
  \par\medskip\Needspace{3\baselineskip}\noindent
  {\bfseries\color{structurecolor}#1}\par\nobreak\smallskip
}{\par\medskip}
\newcommand{\StudyBadge}[2]{%
  \begingroup\small\sffamily
  \ifstrequal{#1}{mid}{\colorlet{studybadge}{green!35!black}}{%
  \ifstrequal{#1}{final}{\colorlet{studybadge}{yellow!45!black}}{\colorlet{studybadge}{red!45!black}}}%
  \textcolor{studybadge}{\rule[-.2ex]{1.2pt}{1.2em}}\enspace\textcolor{black}{#2}\quad\endgroup}
\newcommand{\StudyRef}[3]{%
  \ifcsname studyref@#1@#2\endcsname
    \csname studyref@#1@#2\endcsname{#3}%
  \else #3\textsuperscript{[待定页]}\fi}
\newcommand{\StudyRepoRef}[3]{%
  \href{https://github.com/CrazyShout/CityU-StudyNotes/blob/\StudySourceRef/#1\ifstrempty{#2}{}{\##2}}{#3}}
\newcommand{\StudyNotebookCell}[2]{%
  \ifcsname studynb@#1@#2\endcsname
    \input{\csname studynb@#1@#2\endcsname}%
  \else\PackageError{studynotes}{Missing saved notebook cell #2}{}\fi}
\newcommand{\StudyCoverExtra}{}
\newcommand{\StudyCover}[2]{%
  \thispagestyle{empty}
  \StudyTarget{\StudyBook__root}\StudyCoverExtra
  \vspace*{22mm}
  {\sffamily\large\color{structurecolor}CITYU STUDYNOTES\par}
  \vspace{12mm}
  {\sffamily\bfseries\Huge\color{structurecolor}#1\par}
  \vspace{6mm}{\Large #2\par}
  \vspace{10mm}\noindent\textcolor{structurecolor}{\rule{.42\linewidth}{1pt}}\par
  \vspace{8mm}中文讲解 · English terminology · 双语题答\par
  A4 单栏 · 根据已下载教师材料整理\par
  \vfill
  {\small 本资料用于课程学习与复习，保留原材料、历史题源与教学补充的身份。\par
  模板：ElegantBook 4.7；版式与构建由 CityU-StudyNotes 维护。\par
  \StudyDate\par}
  \clearpage
}
\newcommand{\StudyQuestionExtra}{}
\newcommand{\StudyQuestion}[3]{%
  \Needspace{10\baselineskip}
  \section*{#1 · #2}\phantomsection\addcontentsline{toc}{section}{#1 · #2}%
  \StudyTarget{\StudyBook__#1}%
  \StudyQuestionExtra
  \lowercase{\def\StudyQuestionAnchor{#1}}%
  \noindent\begin{minipage}[t]{.65\linewidth}{\small\sffamily #3}\end{minipage}\hfill
  \begin{minipage}[t]{.33\linewidth}\raggedleft\scriptsize
  \StudyRef{Questions}{\StudyQuestionAnchor}{题目} · \StudyRef{Answers}{\StudyQuestionAnchor}{答案}
  \end{minipage}\par\medskip
}
\InputIfFileExists{\StudyRoot/.build/latex/\StudyBook/refs.tex}{}{}

\renewcommand{\StudySourceRef}{CS5222-Final-2021A}
\renewcommand{\StudyDate}{2026-11-26}

\begin{document}
\StudyCover{CS5222 · 2021/22 Sem A 期末真题双语精析合订本}{全卷 8 道大题（26 小问 · 100 分）· 独立复算、非官方答案勘误与开卷速查}
\tableofcontents

% =============================================================================
\StudyChapter{CS5222-Final-Guide}{2021/22 Sem A 真题结构、开卷合规与参考答案勘误}
% =============================================================================

本册完整整理香港城市大学（CityUHK）\textbf{CS5222 Computer Networks \& Internets（Semester A 2021/22）} 期末考试全部 \textbf{8 道大题（共 26 小问、满分 100 分，时长 2 小时，开卷纸质资料）}。所有配图均为原创 Matplotlib 矢量重绘，\textbf{不含任何教材扫描片段或原卷截图}，符合期末开卷考场规定（允许自带合规计算器、Canvas 教学资料打印件及个人笔记，严禁教材及其打印节选与电子设备）。

\textbf{全部数值均经独立复算核对，并澄清流传非官方参考答案（unofficial docx）的三处细节偏差：}
\begin{enumerate}
\item \textbf{Q7(c) 有限节点 $N=4$ 与无穷极限区分}：流传答案仅给出渐近极限“Slotted ALOHA 最大效率 $1/e\approx 37\%$、Pure ALOHA 为 $1/(2e)\approx 18\%$”，未区分本题给定了 \textbf{$N=4$ 个节点}。精确计算：$N=4$ 时 Slotted ALOHA 在 $p^*=1/4$ 取最大效率 $4(1/4)(3/4)^3 = \mathbf{27/64 = 42.19\%}$，Pure ALOHA 在 $p^*=1/7$ 取最大效率 $4(1/7)(6/7)^6 \approx \mathbf{22.67\%}$；$1/e$ 与 $1/(2e)$ 为 $N\to\infty$ 的渐近极限。作答时将有限 $N=4$ 精确式与渐近极限同时列出最为稳妥。
\item \textbf{Q4(a) 隐含假设说明}：原题未显式指定初始 $\mathrm{ssthresh}$ 数值，默认整个 70 MSS 下载全程处于慢启动阶段（即 $\mathrm{ssthresh}\ge 32\text{ MSS}$）。在答卷第一行写明该前提条件更严谨。
\item \textbf{Q8(d) RTS/CTS 避免的是“长数据帧冲突”}：短 RTS 控制帧之间仍可能发生预约冲突（reservation collision），但一旦接收端广播的 CTS 被覆盖范围内所有潜在干扰站正确接收并设置 NAV，后续长 DATA 帧即可免于隐藏终端冲突。
\end{enumerate}

\begingroup\small\setlength{\tabcolsep}{3pt}\renewcommand{\arraystretch}{1.22}
\begin{longtable}{@{}>{\raggedright\arraybackslash}p{\dimexpr0.10\linewidth-1.50000\tabcolsep\relax}>{\raggedright\arraybackslash}p{\dimexpr0.24\linewidth-1.50000\tabcolsep\relax}>{\raggedright\arraybackslash}p{\dimexpr0.10\linewidth-1.50000\tabcolsep\relax}>{\raggedright\arraybackslash}p{\dimexpr0.56\linewidth-1.50000\tabcolsep\relax}@{}}
\toprule
题号 & 对应章节 & 分值 & 考查要点（对照本学期章节与 Tutorials） \\
\midrule\endfirsthead
题号 & 对应章节 & 分值 & 考查要点（对照本学期章节与 Tutorials） \\
\midrule\endhead
\bottomrule\endfoot
Q1 & Chapter 1 & 14 分 & 传输/传播时延综合对比（公交车运送 vs 拨号）；路由器输出排队判定 \\
Q2 & Ch. 2 / 3 / 4 / 6 & 10 分 & 五项核心协议识别（ARP、DHCP、TCP、CSMA/CD、SMTP） \\
Q3 & Chapter 2 & 11 分 & C/S 与 P2P 最小分发时间下界、P2P 自扩展性与 BitTorrent Tracker \\
Q4 & Chapter 3 & 15 分 & TCP 慢启动下载 RTT 轮数与平均吞吐率；连续报文段序号与确认号推算 \\
Q5 & Chapter 4 & 16 分 & Forwarding 与 Routing 区别；8 位地址最长前缀匹配（LPM）区间计数 \\
\textbf{Q6} & \textbf{Chapter 5} & \textbf{14 分} & \textbf{5 节点图从 $a$ 出发运行 Dijkstra 逐步填表并绘制最短路径树（SPT）} \\
Q7 & Chapter 6 & 8 分 & 4 节点 Slotted ALOHA 静默概率与效率；Pure ALOHA 易损期加倍分析 \\
Q8 & Chapter 7 & 12 分 & 隐藏终端定义、干扰图逐问识别与 802.11 RTS/CTS 信道预约机制 \\
\end{longtable}\endgroup


% =============================================================================
\StudyChapter{CS5222-Final-2021A}{2021/22 Sem A 完整八道真题与双语精解}
% =============================================================================

\StudyHeading{1}{Question 1 · Delay Calculations（14 分｜Chapter 1）}

\StudyHeading{2}{Q1(a)（7 分）｜公交车运送存储介质与 56 kbps 拨号传输比较}
\textbf{Question:} A user needs to transmit $1.5\text{ GBytes}$ of data to a server. Only $56\text{-Kbps}$ dial-up access is available. Alternatively, a bus visits once a day from the city $150\text{ km}$ away (worst-case wait: 1 day), has a $100\text{-Mbps}$ WiFi link to its onboard storage, travels at $60\text{ km/h}$, and uploads to the server over a $1\text{-Gbps}$ link back in the city. Assume $1\text{G}=10^9, 1\text{M}=10^6, 1\text{K}=10^3$. Which way is faster? Explain.

\textbf{中文题译：} 需将 $1.5\text{ GB}$ 数据传至服务器。方案一：56 kbps 拨号直传；方案二：最坏等待每日一班的公交车 1 天（24 h），经 100 Mbps WiFi 拷入车载存储，公交以 60 km/h 驶回 150 km 外的城市，再经 1 Gbps 链路上传服务器。比较两种方案总耗时。

\begin{StudySolution}{标准双语解析 / Verified Answer: Bus is faster (26.54 h vs 59.52 h)}
先统一单位为比特：$L = 1.5 \times 10^9 \text{ Bytes} \times 8 = 1.2 \times 10^{10}\text{ bits}$（$12\text{ Gbits}$）。
\begin{itemize}
\item \textbf{拨号直传（Dial-up）}：
  \[
  T_{\text{dial-up}} = \frac{1.2 \times 10^{10}}{56 \times 10^3} = 214{,}285.71\text{ s} \approx \mathbf{59.524\text{ hours}} \;(2.48\text{ days}). \tag{F.1}
  \]
\item \textbf{公交车运送（Bus）}：最坏等车 $24\text{ h} = 86{,}400\text{ s}$；WiFi 上传 $1.2\times 10^{10}/10^8 = 120\text{ s}$；行车 $150/60 = 2.5\text{ h} = 9{,}000\text{ s}$；城市 1 Gbps 上传 $1.2\times 10^{10}/10^9 = 12\text{ s}$。合计：
  \[
  T_{\text{bus}} = 86{,}400 + 120 + 9{,}000 + 12 = 95{,}532\text{ s} \approx \mathbf{26.537\text{ hours}} \;(1.11\text{ days}). \tag{F.2}
  \]
\end{itemize}
\textbf{结论 / Conclusion:} \textbf{By bus is faster} ($26.54\text{ h} < 59.52\text{ h}$). High-capacity physical transport overcomes a narrowband 56 kbps bottleneck even under the 24-hour worst-case wait.
\end{StudySolution}


\StudyHeading{2}{Q1(b)（7 分）｜两台主机同时发包，路由器 A 会发生排队时延吗？}
\textbf{Question:} Two hosts on the left simultaneously start transmitting 1500-Byte packets at $t=0$ toward Router B via Router A. Both host-to-A links are $4\text{ Mbps}$, with propagation delays of $2\text{ ms}$ and $6\text{ ms}$, respectively. Link A--B is $4\text{ Mbps}$ with $10\text{ ms}$ propagation delay. Will queueing delay occur at Router A? Explain ($1\text{M}=10^6$).

\textbf{中文题译：} 左侧两台主机于 $t=0$ 同时各发 1 个 1500 B 分组经路由器 A 发往 B。两条接入链路均为 4 Mbps、传播时延分别为 2 ms 与 6 ms；A–B 链路亦为 4 Mbps、传播时延 10 ms。路由器 A 会发生排队时延吗？

\StudyImage{CS5222/exam-review/Final/assets/final-2021a-q1b.png}
\StudyCaption{原创绘制 Q1(b) 网络拓扑（左）与路由器 A 收发时间轴（右）：两包在出链路间留有 1 ms 空隙，无排队等待。}

\begin{StudySolution}{标准双语解析 / Verified Answer: No queueing delay}
\begin{enumerate}
\item 单包比特数 $L = 1500 \times 8 = 12{,}000\text{ bits}$；在所有 $4\text{ Mbps}$ 链路上的传输时延均为 $d_{\text{trans}} = 12{,}000 / (4\times 10^6) = \mathbf{3\text{ ms}}$。
\item \textbf{包 1（2 ms 接入链路）}：在 $t = 3 + 2 = \mathbf{5\text{ ms}}$ 完全到达路由器 A（存储转发），立即向出链路 A–B 发送，于 $t = 5 + 3 = \mathbf{8\text{ ms}}$ 完全离开路由器 A。
\item \textbf{包 2（6 ms 接入链路）}：在 $t = 3 + 6 = \mathbf{9\text{ ms}}$ 才完全到达路由器 A。此时出链路早在 $t=8\text{ ms}$ 已恢复空闲。
\end{enumerate}
\textbf{结论 / Conclusion:} \textbf{No queueing delay occurs at Router A}, because Packet 1 finishes leaving A at $t=8\text{ ms}$, strictly before Packet 2 finishes arriving at $t=9\text{ ms}$ (the 10 ms propagation delay on link A--B does not affect link availability at A).
\end{StudySolution}


\StudyHeading{1}{Question 2 · Protocols（10 分｜全课程协议识别）}

\textbf{Question:} Name the protocol supporting each function (abbreviations acceptable):
(a) Translates an IP address into a MAC address in a LAN.
(b) Allows a newly arriving host to obtain an IP address automatically.
(c) A transport-layer protocol that ensures reliability for file transfer.
(d) A MAC protocol used by Ethernet.
(e) An application-layer protocol used to transfer mail between senders' and recipients' mail servers.

\begin{StudySolution}{标准答案 / Verified Answers (2 分 $\times$ 5)}
\begin{itemize}
\item \textbf{(a) ARP}（Address Resolution Protocol，地址解析协议）
\item \textbf{(b) DHCP}（Dynamic Host Configuration Protocol，动态主机配置协议）
\item \textbf{(c) TCP}（Transmission Control Protocol，传输控制协议）
\item \textbf{(d) CSMA/CD}（Carrier Sense Multiple Access with Collision Detection，带冲突检测的载波监听多址接入）
\item \textbf{(e) SMTP}（Simple Mail Transfer Protocol，简单邮件传输协议）
\end{itemize}
\end{StudySolution}


\StudyHeading{1}{Question 3 · Peer-to-Peer (P2P) File Distribution（11 分｜Chapter 2）}

\textbf{Question:} Distribute a file of $F = 10\text{ GBytes}$ to $N = 10$ peers. Server upload rate $u_s = 10\text{ Mbps}$; each peer has download rate $d_i = 2\text{ Mbps}$ and upload rate $u_i = 200\text{ Kbps}$ ($1\text{G}=10^9, 1\text{M}=10^6, 1\text{K}=10^3$).
(a) (3 marks) Minimum client-server distribution time.
(b) (4 marks) Minimum P2P distribution time.
(c) (2 marks) High-level reason why P2P is more scalable than client-server.
(d) (2 marks) How do BitTorrent peers obtain one another's IP addresses?

\begin{StudySolution}{标准双语解析 / Verified Answers}
先换算文件总比特数：$F = 10 \times 10^9 \times 8 = 8 \times 10^{10}\text{ bits}$（$80\text{ Gbits}$）。
\begin{itemize}
\item \textbf{(a) Client-Server 最小分发时间：}
  \[
  \begin{aligned}
  D_{\text{CS}} &= \max\!\left\{\frac{NF}{u_s},\; \frac{F}{d_{\min}}\right\}
  = \max\!\left\{\frac{10\times 8\times 10^{10}}{10^7},\; \frac{8\times 10^{10}}{2\times 10^6}\right\} \\
  &= \max\{80{,}000,\; 40{,}000\} = \mathbf{80{,}000\text{ s}} \;(\approx 22.22\text{ h}).
  \end{aligned}
  \tag{F.3}
  \]
\item \textbf{(b) P2P 最小分发时间：} 系统总上传速率 $u_s + \sum_{i=1}^{10} u_i = 10\text{ Mbps} + 10\times 0.2\text{ Mbps} = 12\text{ Mbps}$：
  \[
  \begin{aligned}
  D_{\text{P2P}} &= \max\!\left\{\frac{F}{u_s},\; \frac{F}{d_{\min}},\; \frac{NF}{u_s + \sum u_i}\right\} \\
  &= \max\!\left\{8{,}000,\; 40{,}000,\; \frac{8\times 10^{11}}{1.2\times 10^7}\right\}
  = \mathbf{66{,}666.67\text{ s}} \;(\approx 18.52\text{ h}).
  \end{aligned}
  \tag{F.4}
  \]
\item \textbf{(c) P2P 高可扩展性原因（Self-scalability）：}
  \textbf{English:} Peers act not only as consumers of bits but also as redistributors—each newly joining peer brings additional upload capacity $u_i$ into the swarm, so aggregate upload capacity scales linearly with $N$, whereas in client-server the fixed server upload rate $u_s$ causes distribution time $NF/u_s$ to grow linearly with $N$.
  \textbf{中文：} P2P 具有\textbf{自扩展性}：每个新加入的对等方既消耗下载带宽，也贡献上传带宽分担分发压力；而 C/S 仅靠单一服务器上传，耗时会随节点数 $N$ 线性增长。
\item \textbf{(d) BitTorrent 获取其他对等方 IP 的机制：}
  新加入 torrent 的节点先联系该种子的追踪服务器 \textbf{Tracker}（登记自身并获得当前活跃对等方的随机子集 IP 列表）；无中心模式下也可通过分布式哈希表 \textbf{DHT} 发现邻居。
\end{itemize}
\end{StudySolution}


\StudyHeading{1}{Question 4 · TCP Congestion Control and Sequence Number（15 分｜Chapter 3）}

\textbf{Question:} TCP starts in Slow Start with $\mathrm{cwnd}=1\text{ MSS}$ and doubles $\mathrm{cwnd}$ each RTT until $\mathrm{ssthresh}$ is reached. A browser downloads a webpage of $70\text{ MSS}$ ($1\text{ MSS} = 1024\text{ Bytes}$).
(a) (3 marks) Assuming no packet loss, how many RTTs are required (including 1 RTT for the TCP handshake)?
(b) (3 marks) If $\mathrm{RTT}=100\text{ ms}$, what is the average throughput in kbps?
(c) (3 marks) In slow start, is throughput higher for a larger webpage than a smaller one? Answer Yes or No and explain.
(d) (6 marks) Host B has received bytes 0–77 from Host A. A sends two back-to-back segments (first: \StudyInline{Seq=78, 12 bytes}; second: \StudyInline{Seq=x, 26 bytes}). Both ACKs (\StudyInline{ACK=y} and \StudyInline{ACK=z}) are lost. Find $x, y, z$.

\StudyImage{CS5222/exam-review/Final/assets/final-2021a-q4d.png}
\StudyCaption{原创绘制 Q4(d) 背靠背 TCP 报文段与丢失确认时序图，对应 $x=90, y=90, z=116$。}

\begin{StudySolution}{标准双语解析 / Verified Answers}
\begin{itemize}
\item \textbf{(a) 所需 RTT 总数：$\mathbf{8\text{ RTTs}}$}（前提假设：全程无丢包且初始 $\mathrm{ssthresh} \ge 32\text{ MSS}$，使下载保持在慢启动指数增长阶段）：
  \begin{itemize}
   \item 第 1 个 RTT：TCP 三次握手（SYN 与 SYN-ACK；第 3 步握手报文捎带发送 HTTP GET 请求）。
   \item 第 2–7 个 RTT（共 6 个数据轮次）：慢启动窗口依次为 $1, 2, 4, 8, 16, 32$，累计传完 $1 + 2 + 4 + 8 + 16 + 32 = \mathbf{63\text{ MSS}}$。
   \item 第 8 个 RTT：窗口可发至 64 MSS，单轮传完剩余 $70 - 63 = \mathbf{7\text{ MSS}}$。
  \end{itemize}
  合计 $1 + 6 + 1 = \mathbf{8\text{ RTTs}}$。
\item \textbf{(b) 平均吞吐率（$\mathrm{RTT}=100\text{ ms}$）：}
  总比特数 $= 70 \times 1024 \times 8 = 573{,}440\text{ bits}$；总耗时 $= 8 \times 0.1 = 0.8\text{ s}$：
  \[
  \mathrm{Throughput} = \frac{573{,}440\text{ bits}}{0.8\text{ s}} = 716{,}800\text{ bps} = \mathbf{716.8\text{ kbps}}. \tag{F.5}
  \]
\item \textbf{(c) 慢启动阶段大网页的平均吞吐率是否高于小网页？—— 是（Yes）。}
  \textbf{English:} \textbf{Yes.} During slow start, $\mathrm{cwnd}$ doubles each RTT ($1,2,4,8,\dots$). A larger webpage transmits later rounds at larger congestion window sizes and amortizes the fixed 1-RTT connection handshake over more bytes, achieving a higher average throughput than a small webpage that finishes while $\mathrm{cwnd}$ is still small.
\item \textbf{(d) 序号与确认号：$x = \mathbf{90},\; y = \mathbf{90},\; z = \mathbf{116}$。}
  第 1 段覆盖字节 $78\text{–}89$（共 12 B），故第 2 段起始字节序号为 $x = 78 + 12 = \mathbf{90}$；两段按序抵达 B，B 收妥第 1 段后期待下一字节 $y = \mathbf{90}$；收妥第 2 段（字节 $90\text{–}115$，共 26 B）后期待下一字节 $z = 90 + 26 = \mathbf{116}$（确认段随后在返回途中丢失，不改变 B 发出它们时填入的确认号）。
\end{itemize}
\end{StudySolution}


\StudyHeading{1}{Question 5 · Forwarding and Routing（16 分｜Chapter 4）}

\textbf{Question:}
(a) (4 marks) What is the difference between forwarding and routing in the network layer?
(b) (12 marks) Consider a datagram network with 8-bit host addresses and longest prefix matching (LPM): Prefix \StudyInline{1} $\to$ Interface 0; \StudyInline{10} $\to$ Interface 1; \StudyInline{111} $\to$ Interface 2; \StudyInline{Otherwise} $\to$ Interface 3. Give the address range and number of addresses for each interface.

\begin{StudySolution}{标准双语解析 / Verified Answers}
\textbf{(a) Forwarding vs. Routing：}
\begin{itemize}
\item \textbf{Forwarding（转发，数据平面）}：单台路由器内部的\textbf{本地动作}，根据本地转发表将入端口到达的分组交换至对应出端口（纳秒级硬件实现）。
\item \textbf{Routing（路由，控制平面）}：全网范围的\textbf{全局过程}，通过路由协议（如 OSPF、BGP）协同计算源到目的的端到端路径，并生成各路由器的转发表（秒级软件实现）。
\end{itemize}

\textbf{(b) 8 位地址最长前缀匹配分配表（共 $2^8 = 256$ 个地址）：}
以 \StudyInline{1} 开头的 128 个地址中，\StudyInline{10xxxxxx}（64 个）匹配更长的 2 位前缀 \StudyInline{10} 走接口 1，\StudyInline{111xxxxx}（32 个）匹配更长的 3 位前缀 \StudyInline{111} 走接口 2，故落入接口 0 的仅剩 \textbf{\StudyInline{110xxxxx}}（32 个地址）。

\begingroup\small\setlength{\tabcolsep}{3pt}\renewcommand{\arraystretch}{1.20}
\begin{longtable}{@{}>{\raggedright\arraybackslash}p{\dimexpr0.14\linewidth-1.50000\tabcolsep\relax}>{\raggedright\arraybackslash}p{\dimexpr0.12\linewidth-1.50000\tabcolsep\relax}>{\raggedright\arraybackslash}p{\dimexpr0.56\linewidth-1.50000\tabcolsep\relax}>{\raggedright\arraybackslash}p{\dimexpr0.18\linewidth-1.50000\tabcolsep\relax}@{}}
\toprule
Prefix & Interface & Range of destination host addresses（二进制 / 十进制） & Number of addresses \\
\midrule\endfirsthead
Prefix & Interface & Range of destination host addresses（二进制 / 十进制） & Number of addresses \\
\midrule\endhead
\bottomrule\endfoot
\texttt{1} & \textbf{0} & \texttt{11000000} -- \texttt{11011111}（十进制 192--223） & \textbf{32} \\
\texttt{10} & \textbf{1} & \texttt{10000000} -- \texttt{10111111}（十进制 128--191） & \textbf{64} \\
\texttt{111} & \textbf{2} & \texttt{11100000} -- \texttt{11111111}（十进制 224--255） & \textbf{32} \\
\texttt{Otherwise} & \textbf{3} & \texttt{00000000} -- \texttt{01111111}（十进制 0--127） & \textbf{128} \\
\end{longtable}\endgroup
总地址数验算：$32 + 64 + 32 + 128 = 256$。
\end{StudySolution}


\StudyHeading{1}{Question 6 · Dijkstra's Algorithm（14 分｜Chapter 5）}

\textbf{Question:} Consider the 5-node weighted network $\{a, b, c, d, e\}$ with edges $(a,b)=8,\; (a,c)=3,\; (a,d)=2,\; (b,d)=1,\; (b,e)=1,\; (c,d)=4,\; (d,e)=5$. Run Dijkstra's shortest-path algorithm from \textbf{node $a$} to all other nodes: (a) (10 marks) Fill in the step-by-step table; (b) (4 marks) Draw the resulting shortest-path tree (SPT) from node $a$.

\StudyImage{CS5222/exam-review/Final/assets/final-2021a-q6.png}
\StudyCaption{原创绘制 Q6 五节点无向带权原图（左）与从源节点 $a$ 出发的最短路径树（右，红实线，括号内为最终最短代价）。}

\begin{StudySolution}{标准双语解析 / Verified Dijkstra Table \& SPT}
\textbf{(a) 逐行标准填表（Step 0–4）：}

\begingroup\small\setlength{\tabcolsep}{3pt}\renewcommand{\arraystretch}{1.20}
\begin{longtable}{@{}>{\raggedright\arraybackslash}p{\dimexpr0.08\linewidth-1.66667\tabcolsep\relax}>{\raggedright\arraybackslash}p{\dimexpr0.20\linewidth-1.66667\tabcolsep\relax}>{\raggedright\arraybackslash}p{\dimexpr0.18\linewidth-1.66667\tabcolsep\relax}>{\raggedright\arraybackslash}p{\dimexpr0.18\linewidth-1.66667\tabcolsep\relax}>{\raggedright\arraybackslash}p{\dimexpr0.18\linewidth-1.66667\tabcolsep\relax}>{\raggedright\arraybackslash}p{\dimexpr0.18\linewidth-1.66667\tabcolsep\relax}@{}}
\toprule
Step & $N'$ & $D(b),p(b)$ & $D(c),p(c)$ & $D(d),p(d)$ & $D(e),p(e)$ \\
\midrule\endfirsthead
Step & $N'$ & $D(b),p(b)$ & $D(c),p(c)$ & $D(d),p(d)$ & $D(e),p(e)$ \\
\midrule\endhead
\bottomrule\endfoot
0 & $a$ & $8, a$ & $3, a$ & \textbf{$2, a$} & $\infty$ \\
1 & $ad$ & \textbf{$3, d$} & $3, a$ & — & $7, d$ \\
2 & $adb$ & — & \textbf{$3, a$} & — & $4, b$ \\
3 & $adbc$ & — & — & — & \textbf{$4, b$} \\
4 & $adbce$ & — & — & — & — \\
\end{longtable}\endgroup

\textbf{逐步松弛过程：}
\begin{itemize}
\item Step 1 选入最小节点 $n^* = d$（$D(d)=2$）：更新 $b$ 为 $\min(8, 2+1)=3,d$；$c$ 候选 $2+4=6>3$ 保持 $3,a$；$e$ 更新为 $\min(\infty, 2+5)=7,d$。
\item Step 2 最小距离 $3$ 在 $b, c$ 间并列（先选 $b$ 或先选 $c$ 最终结果完全一致）：选入 $b$（$D(b)=3$），松弛 $e$ 得 $\min(7, 3+1)=\mathbf{4,b}$。
\item Step 3 选入 $c$（无剩余未定案邻居需更新）；Step 4 选入 $e$，全部节点进入 $N'$，算法终止。
\end{itemize}

\textbf{(b) 最短路径树（SPT，共 $5-1=4$ 条边，见上方右图红实线）：}
按最终前驱 $(p(v), v)$ 连边：\textbf{$(a, d)$（代价 2）、$(a, c)$（代价 3）、$(d, b)$（总代价 3）、$(b, e)$（总代价 4）}。
\end{StudySolution}


\StudyHeading{1}{Question 7 · ALOHA（8 分｜Chapter 6）}

\textbf{Question:} Four active nodes (A, B, C, D) compete using Slotted ALOHA; each transmits in each slot with probability $p$.
(a) (2 marks) Probability that none of the nodes transmit in slot 2.
(b) (2 marks) Efficiency (long-run fraction of successful slots) of this 4-node system.
(c) (4 marks) If Pure ALOHA is used instead of Slotted ALOHA, is the efficiency higher or lower? Why?

\begin{StudySolution}{标准双语解析（含 $N=4$ 精确解与渐近极限）/ Verified Answers}
\begin{itemize}
\item \textbf{(a) 时隙 2 无任何节点发送的概率：} 四个节点相互独立且各时隙独立，全部静默的概率为 $\mathbf{(1-p)^4}$。
\item \textbf{(b) 4 节点 Slotted ALOHA 效率：} 单一时隙恰好 1 个节点发送、其余 3 个节点静默（共 4 种互斥情况）：
  \[
  \eta_{\text{slotted}}(4, p) = \mathbf{4p(1-p)^3}. \tag{F.6}
  \]
\item \textbf{(c) 改用 Pure ALOHA 效率更高还是更低？—— 更低（Lower）。}
  \begin{enumerate}
  \item \textbf{物理原因}：Pure ALOHA 取消了时隙边界同步，任意时刻到达的帧立即发送，使得在 $t_0$ 开始发送的帧会与 $[t_0 - T_{\text{frame}},\; t_0 + T_{\text{frame}}]$ 长度达 \textbf{$2T_{\text{frame}}$} 内启动的任何其他帧发生首尾重叠冲突——\textbf{冲突易损期（vulnerable period）加倍}。
  \item \textbf{精确公式对比（对任意 $0<p<1$）}：4 节点 Pure ALOHA 成功率降为 $\eta_{\text{pure}}(4,p) = 4p(1-p)^{2(4-1)} = 4p(1-p)^6 < 4p(1-p)^3$。
  \item \textbf{精确极值与渐近极限对照}：$N=4$ 时，Slotted ALOHA 在 $p^*=1/4$ 取得峰值 $\mathbf{27/64 \approx 42.19\%}$，Pure ALOHA 在 $p^*=1/7$ 峰值仅 $\mathbf{22.67\%}$；当 $N\to\infty$ 时二者峰值分别收敛到 $1/e \approx 36.8\%$ 与 $1/(2e) \approx 18.4\%$。
  \end{enumerate}
\end{itemize}
\end{StudySolution}


\StudyHeading{1}{Question 8 · Wireless MAC Protocol（12 分｜Chapter 7）}

\textbf{Question:}
(a) (2 marks) What is the ``hidden terminal problem'' in wireless communications? State one undesirable effect it brings.
(b) (3 marks) In the 5-station interference graph below (edges: 1--2, 1--4, 1--5, 2--3, 3--4, 4--5), Station 1 is transmitting to Station 4. Which station(s) can cause the hidden terminal problem?
(c) (3 marks) Suppose Station 1 is transmitting to Station 5. Which station(s) can cause the hidden terminal problem?
(d) (4 marks) State a mechanism in IEEE 802.11 that avoids data frame collisions even in the presence of hidden terminals, and briefly explain how it works.

\StudyImage{CS5222/exam-review/Final/assets/final-2021a-q8.png}
\StudyCaption{原创绘制 Q8 无线干扰关系图及 (b)/(c) 隐藏终端判定结果。}

\begin{StudySolution}{标准双语解析 / Verified Answers}
\begin{itemize}
\item \textbf{(a) 隐藏终端问题定义与不良后果：}
  \textbf{English:} The hidden terminal problem arises when two wireless stations (e.g., A and C) are outside each other's transmission range and cannot sense each other's carrier, yet both lie within range of a common receiver (B). Consequently, both may transmit simultaneously and collide at B. \textbf{Undesirable effect:} Repeated frame collisions at the receiver, causing frame loss, retransmissions, wasted channel airtime, and reduced throughput.
  \textbf{中文：} 两台彼此超出通信半径而互相监听不到的无线站点，同时向共同覆盖范围内的同一接收站发送帧，导致信号在接收端叠加冲突；后果是产生帧丢失、重传开销与吞吐率下降。
\item \textbf{(b) $1 \to 4$ 的隐藏终端：Station 3。}
  使用集合判定式 $\mathrm{Hidden}(s\to r) = (\mathcal{N}(r)\setminus\{s\})\setminus\mathcal{N}(s)$：
  $\mathcal{N}(4)=\{1,3,5\}$，$\mathcal{N}(1)=\{2,4,5\}$ $\Rightarrow$ $\{3,5\}\setminus\{2,4,5\} = \mathbf{\{3\}}$（Station 3 在接收站 4 的覆盖内，但不在发送站 1 的覆盖内；Station 5 能直接听到 1，不是隐藏终端）。
\item \textbf{(c) $1 \to 5$ 的隐藏终端：无（None / No hidden terminal）。}
  $\mathcal{N}(5)=\{1,4\}$，$\mathcal{N}(1)=\{2,4,5\}$ $\Rightarrow$ $\{4\}\setminus\{2,4,5\} = \mathbf{\varnothing}$（接收站 5 除发送站 1 外的唯一邻居是 Station 4，而 Station 4 能直接监听到 Station 1 的载波，不会构成隐藏终端）。
\item \textbf{(d) IEEE 802.11 解决方案：\textbf{RTS/CTS（Request-to-Send / Clear-to-Send）信道预约机制}。}
  \textbf{工作流程（四步）：}
  (1) 发送方监听空闲 DIFS 后发送短 \textbf{RTS} 帧请求预约信道，帧内携带完成本次 DATA+ACK 所需的总时长；
  (2) 接收方（或 AP）在 SIFS 后广播短 \textbf{CTS} 帧确认预约；
  (3) 覆盖范围内所有听到 CTS 的站点（含隐藏终端）读取时长并设置网络分配向量 \textbf{NAV}，在预约期内推迟发送；
  (4) 发送方收到 CTS 后经 SIFS 发送长 \textbf{DATA} 帧，接收方校验无误后经 SIFS 回复 \textbf{ACK}。
  \emph{机制边界：} 短 RTS 帧之间仍可能发生碰撞，但因控制帧很短，开销远小于长 DATA 帧碰撞；一旦 CTS 成功广播，长数据帧即可免受隐藏终端干扰。
\end{itemize}
\end{StudySolution}

\end{document}
